%\SweaveUTF8




\setkeys{Gin}{width = \textwidth}

\begin{Problem}[Stratified simple random sampling]
    Use the file \code{sepe.dat} again. Define the stratum indicator \code{strat
    = 1} for establishments with \code{empl} $\leq 250$ and \code{strat = 2}
    otherwise. Also define \code{Nh}, which gives the population size of the
    corresponding stratum for each establishment. Add both variables to the
    data set.

    Draw a stratified simple random sample without replacement, with
    $n_1 = 100$ establishments from the first stratum and $n_2 = 50$ from the
    second stratum. Estimate the population mean of environmental expenditures
    and its standard error.

    Finally, compare the estimated standard error with that from a simple
    random sample of $n = 150$ establishments. For the simple random sample,
    compute the standard error both with and without the finite population
    correction.

    \begin{Rhint}
        Load the data \code{sepe.dat} data and define \code{strat} and
        \code{Nh} as described above. Attach the \code{survey} package to the
        search path and draw the stratified sample:


\begin{Scode}
library("survey")
set.seed(11)
sepe_sample <- sepe[stratsample(sepe$strat, c("1" = 100, "2" = 50)), ]
\end{Scode}

        The function \code{stratsample()} is provided by the \code{survey}
        package. To reproduce the solution, use \code{set.seed(11)} before
        drawing the sample. Next, define the survey design object by

\begin{Scode}
dn_sepe <- svydesign(ids = ~1, strata = ~strat, fpc = ~Nh,
                     data = sepe_sample)
\end{Scode}

        Then use \code{svymean()} to estimate the population mean and
        \code{SE()} to extract its estimated standard error.
    \end{Rhint}

\begin{Solution}

\begin{Schunk}
\begin{Sinput}
> sepe <- read.table("http://stat.ethz.ch/Teaching/Datasets/WBL/sepe.dat",
                     header = TRUE)
\end{Sinput}
\end{Schunk}

    Define the strata and their population sizes:

\begin{Schunk}
\begin{Sinput}
> sepe$strat <- ifelse(sepe$empl <= 250, 1, 2)
> Nh <- table(sepe$strat)
> sepe$Nh <- as.numeric(Nh[as.character(sepe$strat)])
> Nh
\end{Sinput}
\begin{Soutput}
  1   2 
580 220 
\end{Soutput}
\end{Schunk}

    Draw the stratified sample:

\begin{Schunk}
\begin{Sinput}
> library("survey")
> set.seed(11)
> idx <- stratsample(sepe$strat, c("1" = 100, "2" = 50))
> sepe_sample <- sepe[idx, ]
> table(sepe_sample$strat)
\end{Sinput}
\begin{Soutput}
  1   2 
100  50 
\end{Soutput}
\end{Schunk}

    Define the design and estimate the population mean:

\begin{Schunk}
\begin{Sinput}
> dn_sepe <- svydesign(ids = ~1, strata = ~strat, fpc = ~Nh,
                       data = sepe_sample)
> m_y <- svymean(~expend, dn_sepe)
> m_y
\end{Sinput}
\begin{Soutput}
        mean     SE
expend 86472 4019.5
\end{Soutput}
\begin{Sinput}
> SE(m_y)
\end{Sinput}
\begin{Soutput}
         expend
expend 4019.504
\end{Soutput}
\end{Schunk}

    Thus, the estimated standard error of the stratified estimator is
    4'020.

%FIXME: use survey; see bottom; above: use dn_stsi instead of dn_sepe


    \begin{Remark}
        For comparison, we draw a separate SRS of size $150$:

\begin{Schunk}
\begin{Sinput}
> set.seed(11)
> srs_sample <- sepe[sample(seq_len(nrow(sepe)), size = 150), ]
\end{Sinput}
\end{Schunk}

        We add population size to \code{srs\_sample} and define the design
        object:

\begin{Schunk}
\begin{Sinput}
> srs_sample$N <- 800
> dn_srs <- svydesign(ids = ~1, fpc = ~N, data = srs_sample)
\end{Sinput}
\end{Schunk}

        Ror the SRS sample, the estimated std. err. of the mean is

\begin{Schunk}
\begin{Sinput}
> SE(svymean(~expend, dn_srs))
\end{Sinput}
\begin{Soutput}
         expend
expend 4729.543
\end{Soutput}
\end{Schunk}

        which is considerably large than the estimated std. err. in case of the
        stratified sample, which is (approximately)
        4'020.

        \vspace{0.5em}

%         Note: A stratified sample can also be drawn using base \textsf{R} only.
%         For example:

% <<>>=
% set.seed(11)
% idx <- unlist(mapply(
%     sample,
%     split(seq_len(nrow(sepe)), sepe$strat),
%     c(100, 50),
%     SIMPLIFY = FALSE
% ))
% @
    \end{Remark}
\end{Solution}
\end{Problem}
